% paper_fci_molecules.tex — GeoVac LCAO FCI Paper (Molecules)
% Topological Full Configuration Interaction for Heteronuclear Diatomics
% Status: Revised, March 12, 2026 — corrected for Hamiltonian balance
% Target: arXiv (physics.comp-ph / quant-ph)
\documentclass[aps,pra,twocolumn,superscriptaddress,longbibliography]{revtex4-2}

\usepackage{amsmath,amssymb,amsthm,graphicx,xcolor,bm,braket,dcolumn,booktabs}
\usepackage{hyperref}

\newtheorem{theorem}{Theorem}
\newtheorem{corollary}{Corollary}

\begin{document}

\title{Topological Full Configuration Interaction for Heteronuclear\\
Diatomics: LiH as a Benchmark}
\thanks{GeoVac FCI Paper in the Geometric Vacuum series}

\author{J.~Loutey}
\affiliation{Independent Researcher}

\date{March 12, 2026}

%% ========================================================================
%% ABSTRACT
%% ========================================================================

\begin{abstract}
We present a Topological LCAO architecture for heteronuclear full
configuration interaction (FCI) within the GeoVac discrete graph
framework.  Two independent atom-centered hydrogenic lattices are
concatenated into a bipartite molecular graph via Mulliken cross-nuclear
attraction on the diagonal, Slater--Condon two-electron repulsion
integrals with exact $R^k$ radial integrals and Gaunt coefficients, and
STO-overlap bridges on the off-diagonal; the resulting one- and
two-electron Hamiltonian is diagonalized in the full Slater determinant
space.  For LiH at $n_{\max}=3$ (28 spatial orbitals, 56 spin-orbitals,
367\,290 determinants), the Boys--Bernardi counterpoise-corrected vertical
binding energy at the experimental geometry ($R=3.015$~bohr) is
$D_{e}^{\mathrm{CP}}=0.110$~Ha, within 20\% of the experimental value
$D_{e}=0.092$~Ha.  The balanced Hamiltonian produces a monotonically
attractive potential energy surface with no equilibrium geometry, a
limitation traced to the $R$-independence of the graph Laplacian's
intra-atom kinetic energy: unlike continuous quantum mechanics, the
discrete adjacency within each atomic subgraph does not respond to
inter-nuclear compression.  A Fock-weighted kinetic correction, derived
from the energy-shell constraint $p_0^2 = Z^2/n^2$ of Paper~7, partially
restores the equilibrium by concentrating orthogonalization repulsion on
compact core orbitals, achieving $R_{\mathrm{eq}}\approx 3.8$~bohr and
$D_e^{\mathrm{CP}}\approx 0.046$~Ha.  A 29-version diagnostic arc
(v0.9.8--v0.9.37) systematically tested and eliminated every competing
hypothesis: basis set superposition error as a geometric driver, integral
incompleteness, orbital exponent relaxation, SO(4) Sturmian bond sphere
geometries, and Hamiltonian imbalance between cross-nuclear and
electron repulsion orbital subspaces.  In particular, Wigner $D$-matrix
selection rules $D^{2}_{10,10}(\gamma)=1$ and
$D^{2}_{00,10}(\gamma)=0$ are proven as structural results.  The
excitation-driven Direct CI solver achieves
$O(N_{\mathrm{SD}}\times N_{\mathrm{connected}})$ scaling, preserving the
$O(V)$ character of the underlying graph Hamiltonian.
\end{abstract}

\maketitle

%% ========================================================================
\section{Introduction}
\label{sec:intro}
%% ========================================================================

The GeoVac program posits that continuous quantum mechanics emerges from
discrete graph topologies: a scale-invariant graph Laplacian defined on
the unit three-sphere $S^3$ encodes the full single-particle
Schr\"odinger equation without any continuous spatial
integration~\cite{paper0,paper1,paper7}.  Prior work established this
correspondence rigorously for single-center atomic systems, recovering
the hydrogenic spectrum as a combinatorial eigenvalue problem with $O(V)$
scaling.  Extending this framework to multi-center molecules, however,
presents a fundamental geometric challenge: two or more atomic lattices,
each living on its own $S^3$ conformal manifold, must be bridged into a
single molecular graph without resorting to the continuous
$\mathbb{R}^3$ integration that the discrete premise is designed to
replace.

In this work we introduce the \emph{Topological LCAO} architecture,
which solves the multi-center problem by concatenating atomic subgraphs
into a bipartite molecular graph connected by discrete off-diagonal
bridges.  The resulting block Hamiltonian is fed into an $O(V)$
excitation-driven Full Configuration Interaction solver that operates
entirely in Slater determinant space.  As a headline result, the
Boys--Bernardi counterpoise-corrected vertical binding energy of LiH at
$n_{\max}=3$ and $R=3.015$~bohr is
$D_{e}^{\mathrm{CP}}=0.110$~Ha, within 20\% of the experimental value
of $0.092$~Ha~\cite{huber}.  The balanced Hamiltonian reveals a
fundamental limitation of the discrete LCAO framework: the absence of
$R$-dependent kinetic repulsion, which is partially resolved by a
Fock-weighted kinetic correction derived from Paper~7~\cite{paper7}.

LiH is the minimal four-electron heteronuclear diatomic and therefore
the ideal first benchmark for the Topological LCAO framework.  Its four
electrons permit exact FCI treatment without truncation of the many-body
expansion, while the $Z_A=3$, $Z_B=1$ asymmetry breaks the inversion
symmetry exploited by homonuclear codes and exposes every heteronuclear
coupling pathway.  The experimental targets are well-established:
$D_e = 0.092$~Ha and $R_{\mathrm{eq}}=3.015$~bohr~\cite{huber}.  This
paper formally extends the companion computational framework described
in Paper~8~\cite{paper8}.

The paper is organized as follows.  Section~\ref{sec:architecture}
develops the bipartite graph construction and the $O(V)$ Direct CI
solver.  Section~\ref{sec:results} presents the LiH benchmark results,
including the counterpoise-corrected potential energy surface and the
Hamiltonian balance requirement.
Section~\ref{sec:diagnostic} documents the 29-version diagnostic arc
that systematically falsified every hypothesis for the observed
behavior---integral incompleteness, orbital
relaxation, and Sturmian geometries.
Section~\ref{sec:kinetic} identifies the root cause as missing
$R$-dependent kinetic repulsion and presents a Fock-weighted correction
derived from Paper~7.
Section~\ref{sec:conclusion} summarizes and outlines next steps.

%% ========================================================================
\section{Topological LCAO Architecture}
\label{sec:architecture}
%% ========================================================================

%% ------------------------------------------------------------------------
\subsection{Subgraph Concatenation}
\label{sec:subgraph}
%% ------------------------------------------------------------------------

The \texttt{MolecularLatticeIndex} instantiates two independent atomic
subgraphs: a lithium lattice with nuclear charge $Z_A=3$ and a hydrogen
lattice with $Z_B=1$.  Each subgraph is a scale-invariant hydrogenic
lattice defined on its own $S^3$ conformal manifold~\cite{paper7},
retaining its atom-centered conformal structure and the exact eigenvalue
spectrum $\varepsilon_{n} = -Z^2/(2n^2)$.  At $n_{\max}=3$, the lithium
subgraph contributes $\sum_{n=1}^{3}n^2 = 14$ spatial orbitals and the
hydrogen subgraph contributes another 14, yielding a combined molecular
basis of 28 spatial orbitals (56 spin-orbitals).

The concatenation is a direct sum at the graph level: the adjacency
matrix of the molecular graph is block-diagonal in the two atomic
subgraphs prior to bridging.  No spatial overlap integral is computed
between the subgraphs at this stage; the topological identity of each
atomic lattice is preserved exactly.

%% ------------------------------------------------------------------------
\subsection{Topological Bridging and the Hamiltonian}
\label{sec:bridging}
%% ------------------------------------------------------------------------

Bonding emerges from a rank-update concatenation that introduces
off-diagonal coupling between the two atomic blocks.  The one-electron
molecular Hamiltonian takes the block form
\begin{equation}
\mathbf{H}_1 =
\begin{pmatrix}
\mathbf{H}_A & \mathbf{V}_{AB} \\
\mathbf{V}_{AB}^{\dagger} & \mathbf{H}_B
\end{pmatrix},
\label{eq:block_h1}
\end{equation}
where $\mathbf{H}_A$ and $\mathbf{H}_B$ are the isolated atomic
Hamiltonians and $\mathbf{V}_{AB}$ encodes the inter-atomic coupling.

The diagonal blocks carry exact graph Laplacian eigenvalues
$\varepsilon_n = -Z^2/(2n^2)$ supplemented by the Mulliken
cross-nuclear attraction: each orbital $\phi_i^{(A)}$ centered on atom
$A$ experiences the potential of nucleus $B$ via the diagonal matrix
element
\begin{equation}
\langle \phi_i^{(A)} | {-Z_B}/{r_B} | \phi_i^{(A)} \rangle,
\label{eq:mulliken}
\end{equation}
evaluated analytically from the hydrogenic Fourier representation.
Off-diagonal elements of $\mathbf{V}_{AB}$ are STO-overlap bridges
weighted by conformal factors derived from the stereographic
projection~\cite{fock1935}.

The two-electron repulsion $V_{ee}$ is constructed via Slater--Condon
rules with two classes of integrals.  Same-atom Coulomb and exchange
integrals use exact Slater $R^k$ radial integrals with Gaunt angular
coefficients, validated against the known analytical values
$F^0(1s,1s)=5Z/8$, $F^0(1s,2s)=17Z/81$, and
$F^0(2s,2s)=77Z/512$~\cite{shull_lowdin}.  Cross-atom Coulomb repulsion
is evaluated via multipole expansion, with exchange restricted to
$s$-orbital pairs where the overlap is non-negligible.

%% ------------------------------------------------------------------------
\subsection{\texorpdfstring{$O(V)$}{O(V)} Direct Configuration Interaction}
\label{sec:direct_ci}
%% ------------------------------------------------------------------------

Distributing four electrons across 56 spin-orbitals generates
$\binom{56}{4} = 367\,290$ Slater determinants.  A na\"ive pairwise
construction of the CI Hamiltonian matrix would require
$O(N_{\mathrm{SD}}^2)$ evaluations---over $10^{11}$ matrix element
computations---rendering exact diagonalization impractical.

We instead adopt the excitation-driven Direct CI algorithm of Knowles
and Handy~\cite{knowles_handy}, which constructs the Hamiltonian action
$\sigma = \mathbf{H}|\Psi\rangle$ by traversing only the nonzero edges
generated by valid single and double excitations from each determinant.
The computational cost scales as
$O(N_{\mathrm{SD}} \times N_{\mathrm{connected}})$, where
$N_{\mathrm{connected}}$ is the average number of determinants connected
to each reference by at most a double excitation.  For the LiH basis at
$n_{\max}=3$, this reduces wall time from hours to approximately
8~seconds when the Slater--Condon assembly loops are JIT-compiled
via Numba~\cite{numba} (approximately 75~seconds in pure Python).

The one-electron integrals $h_{pq}$ and two-electron integrals
$(pq|rs)$ are stored as dense arrays (56 spin-orbitals yield lookup
tables of order $56^2$ and $28^4$, both fitting comfortably in memory),
eliminating the $\sim\!24\,\mu$s per-access overhead of
\texttt{scipy.sparse} index validation that dominated earlier
implementations.

%% ------------------------------------------------------------------------
\subsection{Counterpoise Correction Framework}
\label{sec:cp}
%% ------------------------------------------------------------------------

Basis set superposition error (BSSE) is an inherent artifact of any
atom-centered expansion: at finite basis size, each atom can
variationally borrow orbitals from its partner, artificially lowering
the molecular energy below the true complete-basis-set limit.  We
implement the Boys--Bernardi counterpoise (CP)
correction~\cite{boys_bernardi} natively within the Topological LCAO
framework.

Ghost fragments are constructed by preserving the full bipartite
topology of the molecular graph while zeroing the nuclear charge and
electron population of the partner center.  Specifically, a ghost-$B$
calculation retains all 28 spatial orbitals but sets $Z_B=0$,
eliminating all bridges, cross-nuclear attraction, and cross-atom
$V_{ee}$ involving center $B$.  The fragment energy
$E_A^{\mathrm{mol}}$ is then computed in the identical topological basis
space as the supermolecule, guaranteeing a consistent subtraction:
\begin{equation}
D_e^{\mathrm{CP}} = E_A^{\mathrm{mol}} + E_B^{\mathrm{mol}} -
E_{AB}(R_{\mathrm{eq}}).
\label{eq:cp}
\end{equation}

%% ========================================================================
\section{LiH Benchmark Results}
\label{sec:results}
%% ========================================================================

%% ------------------------------------------------------------------------
\subsection{Computational Parameters}
\label{sec:params}
%% ------------------------------------------------------------------------

The molecular graph is assembled from two atomic subgraphs at
$n_{\max}=3$: a lithium lattice ($Z_A=3$, 14 spatial orbitals) and a
hydrogen lattice ($Z_B=1$, 14 spatial orbitals), yielding a combined
basis of 28 spatial orbitals and 56 spin-orbitals.  The full
four-electron Slater determinant space contains $\binom{56}{4} =
367\,290$ configurations.  The exact ground-state eigenvalue is obtained
from the excitation-driven Direct CI solver described in
Sec.~\ref{sec:direct_ci}, with wall time of approximately 8~seconds on
a single CPU core (Numba-accelerated assembly; approximately 75~seconds
without JIT compilation).

%% ------------------------------------------------------------------------
\subsection{Energetic Accuracy and Basis Set Superposition}
\label{sec:energy}
%% ------------------------------------------------------------------------

Table~\ref{tab:lih_energies} summarizes the key energetic results.  The
raw (uncorrected) binding energy significantly overestimates experiment,
a clear signature of BSSE in the finite atom-centered basis.  The
Boys--Bernardi counterpoise correction quantifies the superposition
error at $\Delta E_{\mathrm{BSSE}} = 0.115$~Ha, with the lithium
fragment contributing $0.105$~Ha and the hydrogen fragment
$0.010$~Ha---reflecting the far greater variational flexibility
available to the smaller atom when it can access the partner's
higher-$Z$ orbitals.

After CP correction, the vertical binding energy at $R=3.015$~bohr is
$D_e^{\mathrm{CP}} = 0.110$~Ha, a 20\% overestimate of the experimental
value $0.092$~Ha~\cite{huber}.  We emphasize that this is a
\emph{vertical} binding energy computed at fixed internuclear distance,
not an adiabatic value at the theoretical equilibrium: the balanced
Hamiltonian produces no equilibrium geometry (Sec.~\ref{sec:contraction}).

\begin{table}[h]
\centering
\begin{tabular}{lc}
\toprule
Quantity & Value \\
\midrule
$n_{\max}$ & 3 \\
Spatial orbitals & 28 \\
Spin-orbitals & 56 \\
Slater determinants & 367\,290 \\
$E_{\mathrm{LiH}}(R{=}3.015)$ (Ha) & $-8.118$ \\
$E_{\mathrm{Li}}^{\mathrm{mol}}$ (Ha) & $-7.392$ \\
$E_{\mathrm{H}}^{\mathrm{mol}}$ (Ha) & $-0.500$ \\
$\Delta E_{\mathrm{BSSE}}$ (Ha) & $-0.115$ \\
$D_e^{\mathrm{raw}}$ (Ha) & 0.226 \\
$D_e^{\mathrm{CP}}$ (Ha) & 0.110 \\
$D_e^{\mathrm{expt}}$ (Ha) & 0.092 \\
Error (\%) & 20 \\
\bottomrule
\end{tabular}
\caption{LiH energetics at $n_{\max}=3$, $R=3.015$~bohr (balanced
\texttt{exact+True} configuration).  The Boys--Bernardi counterpoise
correction removes 0.115~Ha of basis set superposition error.  The
20\% overestimate of $D_e$ is a vertical binding energy at fixed $R$;
the balanced Hamiltonian produces no equilibrium geometry.}
\label{tab:lih_energies}
\end{table}

\paragraph{Hamiltonian balance requirement.}
The cross-nuclear attraction and electron repulsion must operate on
matched orbital subspaces.  Two balanced configurations---\texttt{exact}
cross-nuclear with full cross-atom $V_{ee}$ (\texttt{exact+True}), and
\texttt{fourier} cross-nuclear with $s$-only $V_{ee}$
(\texttt{fourier+s\_only})---produce molecular energies agreeing within
2~mHa across the entire PES, confirming variational consistency.  An
\emph{imbalanced} configuration (\texttt{exact+s\_only}) artificially
overbinds by 0.55~Ha, as the $p$ and $d$ orbitals receive cross-nuclear
attraction without the compensating electron repulsion screening.  An
earlier version of this work inadvertently used the imbalanced
configuration; the results presented here are from the balanced
\texttt{exact+True} Hamiltonian.

%% ------------------------------------------------------------------------
\subsection{Absence of Equilibrium Geometry}
\label{sec:contraction}
%% ------------------------------------------------------------------------

The balanced configurations (\texttt{exact+True} and
\texttt{fourier+s\_only}) produce identical results within 2~mHa,
confirming variational consistency.  However, both yield monotonically
attractive PES curves with no equilibrium geometry.  The equilibrium
reported in earlier versions arose from an imbalanced Hamiltonian
configuration that has since been corrected.

The CP-corrected binding energy $D_e^{\mathrm{CP}}$ decreases smoothly
from 0.464~Ha at $R=1.5$~bohr to 0.042~Ha at $R=5.0$~bohr, with the
PES slope diminishing but never reversing sign.  At no point does the
repulsive kinetic wall balance the attractive cross-nuclear and exchange
interactions: the system is variationally attracted to short bond
lengths at all $R$.

This monotonic behavior is traced to a structural property of the
Topological LCAO: the graph Laplacian kinetic energy within each atomic
subgraph is \emph{completely $R$-independent}.  In continuous quantum
mechanics, bringing two atoms together forces their electrons into
higher-momentum orthogonal states, creating the kinetic repulsion that
establishes the equilibrium bond length.  In the discrete framework,
each atom's adjacency matrix---and hence its contribution to the kinetic
energy via the degree matrix $D$---is fixed at construction time and
does not respond to the approach of a second center.

The only $R$-dependent kinetic contribution comes from the bridge edges
connecting the two subgraphs.  However, these bridge elements are
approximately two orders of magnitude too weak: the maximum bridge
matrix element at $R=3.015$~bohr is 0.017~Ha, compared to the $\sim$0.5~Ha
kinetic correction required by virial theorem analysis (the observed
virial ratio $\eta = -V/(2T) \approx 30$ instead of the equilibrium
value $\eta = 1$).

This absence of equilibrium geometry motivates both the systematic
diagnostic arc in the following section and the kinetic repulsion
analysis in Sec.~\ref{sec:kinetic}.

%% ========================================================================
\section{Systematic Diagnostic Arc}
\label{sec:diagnostic}
%% ========================================================================

The initial implementation (v0.9.8) exhibited heavy overbinding
($D_e^{\mathrm{raw}} = 0.205$~Ha, more than twice the experimental
value) and a 17\% geometric contraction of the equilibrium bond length.
Over the course of 29 controlled experiments spanning versions v0.9.8
through v0.9.36, every plausible hypothesis for these discrepancies was
systematically tested.  The experiments are grouped below by hypothesis
class; the complete version history is tabulated in
Appendix~\ref{app:versions}.

%% ------------------------------------------------------------------------
\subsection{Hypothesis 1: Basis Set Superposition Error
(v0.9.8--v0.9.11)}
\label{sec:hyp_bsse}
%% ------------------------------------------------------------------------

The first hypothesis attributed both the energetic overestimation and
the geometric contraction to BSSE: electrons borrowing orbitals from the
partner center would artificially lower the molecular energy, inflating
$D_e$ and potentially distorting the PES minimum.

The Boys--Bernardi counterpoise correction (Sec.~\ref{sec:cp})
quantified the superposition error at
$\Delta E_{\mathrm{BSSE}} = 0.115$~Ha and corrected the vertical
binding energy to $D_e^{\mathrm{CP}}=0.110$~Ha (20\% error).  BSSE
is the dominant source of energetic error in the raw binding energy
but does not address the absence of an equilibrium geometry, which
has a distinct root cause (Sec.~\ref{sec:kinetic}).

%% ------------------------------------------------------------------------
\subsection{Hypothesis 2: Integral Incompleteness
(v0.9.12--v0.9.13, v0.9.35)}
\label{sec:hyp_integrals}
%% ------------------------------------------------------------------------

The second hypothesis class attributed the geometric contraction to
missing or approximate molecular integrals: either the Mulliken
cross-nuclear attraction overestimates the short-range potential, or the
restriction of cross-atom $V_{ee}$ to $s$-orbital pairs neglects
significant contributions from higher angular momentum channels.

Exact cross-nuclear attraction matrix elements, computed analytically
from the hydrogenic Fourier representation and validated against the
shell theorem, deviate by less than $0.6$~mHa from the Mulliken
diagonal approximation, confirming that the Mulliken scheme is
essentially exact for this system.

Extending the cross-atom $V_{ee}$ to all $(n,l)$ pairs---392 Coulomb
integrals at $n_{\max}=3$ compared to 18 for the $s$-only
restriction---contributes a correction of only $+0.449$~mHa to the
total energy.  This negligible shift reflects the fact that the LiH
ground state is overwhelmingly $s$-orbital dominated: the $1s^2 2s^2$
configuration of lithium and the $1s$ of hydrogen leave higher angular
momentum channels essentially unoccupied.

Integral incompleteness is conclusively falsified as a mechanism for the
geometric contraction.

%% ------------------------------------------------------------------------
\subsection{Hypothesis 3: Orbital Exponent Relaxation (v0.9.36)}
\label{sec:hyp_relaxation}
%% ------------------------------------------------------------------------

The third hypothesis proposed that fixing atomic orbital exponents at
their free-atom values forces artificial proximity: if electrons cannot
polarize toward the bonding region by adjusting their radial extent, the
system compensates by contracting $R_{\mathrm{eq}}$.

A variational orbital exponent parameter $\zeta_B$ was introduced for
the hydrogen center, allowing the effective nuclear charge
$Z_{\mathrm{eff}} = \zeta_B Z_B$ to optimize the radial extent of all
hydrogen-centered orbitals.  The variational minimum occurs at
$\zeta_B \approx 1.30$ at $R=3.015$~bohr, relaxing toward unity at
large $R$ as expected.

The relaxation energy is approximately $-0.27$~Ha, but it is
\emph{uniform across the entire bonding region}: the energy lowering at
$R=2.5$~bohr and $R=3.015$~bohr differs by only $+0.007$~Ha.
Consequently, $\Delta R_{\mathrm{eq}} = 0.0$~bohr---the PES minimum
does not shift.  Moreover, the CP-corrected binding energy inflates by a
factor of $4.7\times$ to $D_e^{\mathrm{CP}}=0.433$~Ha, demonstrating
that the variational freedom is exploited for unphysical basis
compression rather than physical polarization.

Orbital exponent relaxation is eliminated as a mechanism for geometric
contraction.

%% ------------------------------------------------------------------------
\subsection{Hypothesis 4: Bond Sphere and Sturmian Basis
(v0.9.15--v0.9.34)}
\label{sec:hyp_sturmian}
%% ------------------------------------------------------------------------

The fourth and most extensive hypothesis class proposed that the
flat-space LCAO concatenation is fundamentally incompatible with the
$S^3$ topology, and that a proper molecular treatment requires Sturmian
basis functions defined on a single \emph{bond sphere}---an $S^3$
manifold with two weighted poles---connected by SO(4) Wigner $D$-matrix
elements.

This program was pursued across 20 versions (v0.9.15--v0.9.34),
progressing from the initial SO(4) Wigner $D$-matrix implementation (44
tests, v0.9.15) through geometric-mean coupling (v0.9.16, unbound),
Shibuya--Wulfman off-diagonal nuclear attraction (v0.9.17, unbound),
hybrid Fourier--SW schemes (v0.9.18, 55\% overbinding), atom-centered
Sturmian variants (v0.9.20--v0.9.28, all unbound or massively
overbinding), prolate spheroidal MO projections
(v0.9.29--v0.9.30, unphysical), full MO Sturmian FCI
(v0.9.31--v0.9.33, 12--25\% error), and dual-$p_0$ extensions (v0.9.34,
25\% error with $p_0$ values collapsing to a single geometric mean).

The definitive negative result is the structural theorem proven in
companion Paper~8~\cite{paper8} (Theorem~3): in any Sturmian basis
sharing a single momentum parameter $p_0$, the Hamiltonian matrix is
proportional to the overlap matrix,
$\mathbf{H} \propto \mathbf{S}$.  The generalized eigenvalues are
therefore independent of the bond angle $\gamma$, yielding pure $1/R$
nuclear repulsion with no electronic binding.  Recovering
$R$-dependent binding energies requires solving the continuous
two-center prolate spheroidal PDE for the spectrum $\beta_k(R)$,
reintroducing exactly the continuous spatial geometry that the discrete
graph framework is designed to eliminate.  The Bond Sphere Sturmian
alternative is mathematically falsified as a viable discrete framework.

\medskip

Having falsified BSSE as a geometric driver, eliminated integral
incompleteness, ruled out orbital exponent relaxation, and
mathematically disproven the Sturmian alternative, the $O(V)$
Topological LCAO remains the only sound discrete framework for
heteronuclear FCI.  The absence of an equilibrium geometry in the
balanced Hamiltonian is traced to the $R$-independence of the graph
Laplacian kinetic energy, analyzed in the following section.

%% ========================================================================
\section{Missing Kinetic Repulsion and the Fock-Weighted Correction}
\label{sec:kinetic}
%% ========================================================================

%% ------------------------------------------------------------------------
\subsection{The Kinetic Diagnostic}
\label{sec:kinetic_diagnostic}
%% ------------------------------------------------------------------------

The virial theorem for a Coulombic system requires
$\eta \equiv -V/(2T) = 1$ at the equilibrium geometry.
Table~\ref{tab:virial} documents the virial ratio across the LiH PES:
$\eta$ ranges from $46$ at $R=1.5$~bohr to $22$ at $R=5.0$~bohr,
indicating that the kinetic energy is $\sim$30$\times$ too weak relative
to the potential energy at all internuclear distances.

\begin{table}[h]
\centering
\begin{tabular}{ccccc}
\toprule
$R$ (bohr) & $T$ (Ha) & $V$ (Ha) & $\eta=-V/(2T)$ & $E$ (Ha) \\
\midrule
1.5 & $-0.092$ & $-8.380$ & 45.8 & $-8.471$ \\
2.5 & $-0.122$ & $-8.055$ & 32.9 & $-8.178$ \\
3.015 & $-0.133$ & $-7.985$ & 30.0 & $-8.118$ \\
4.0 & $-0.172$ & $-7.908$ & 23.0 & $-8.079$ \\
5.0 & $-0.181$ & $-7.867$ & 21.7 & $-8.049$ \\
\bottomrule
\end{tabular}
\caption{Virial analysis of the balanced LiH PES.  The virial ratio
$\eta$ should equal 1 at equilibrium but ranges from 22 to 46,
indicating the kinetic energy is $\sim$30$\times$ too weak.}
\label{tab:virial}
\end{table}

Direct examination confirms that the intra-atom kinetic energy (the
off-diagonal $\mathbf{H}_1$ within each atomic subgraph) is identically
$R$-independent: the lithium and hydrogen graph Laplacians are
constructed once from the isolated-atom topology and never modified as
$R$ changes.  Bridge edges contribute a small $R$-dependent kinetic
term through the degree matrix, but the maximum bridge element at
$R=3.015$~bohr is only 0.017~Ha, two orders of magnitude below the
required repulsion.

This is a structural limitation of the Topological LCAO architecture:
the discrete graph Laplacian has no mechanism to represent the
orthogonalization kinetic energy that arises when atomic orbitals
overlap in continuous quantum mechanics.

%% ------------------------------------------------------------------------
\subsection{The Fock-Weighted Kinetic Correction}
\label{sec:fock_correction}
%% ------------------------------------------------------------------------

The energy-shell constraint of Paper~7---$p_0^2 = -2E = Z^2/n^2$ for
the $n$-th shell---provides a natural scale for the
``information density'' of each orbital on the conformal $S^3$ manifold.
We exploit this to construct a physically motivated kinetic correction.

For each orbital $\phi_a$ on atom $A$, the Fock-weighted overlap
kinetic correction adds to the diagonal of $\mathbf{H}_1$:
\begin{equation}
\Delta h_{aa} = \lambda \sum_{b \in B}
S_{ab}^2 \left(\frac{Z_A}{n_a}\right)^2
\left(\frac{Z_B}{n_b}\right)^2,
\label{eq:fock_correction}
\end{equation}
where $S_{ab} = \langle \phi_a | \phi_b \rangle$ is the STO overlap
integral between $s$-orbitals on different centers, and $\lambda$ is a
single adjustable parameter.  The Fock weights suppress diffuse
orbitals (high $n$) and concentrate the kinetic penalty on compact core
orbitals whose overlap decays rapidly with $R$.

The physical content of Eq.~(\ref{eq:fock_correction}) is that the cost
of overlapping two topological information structures scales with the
product of their information densities $p_0^2 = Z^2/n^2$.  Core--core
pairs (Li~1$s$--H~1$s$) carry weight $9 \times 1 = 9$, while
diffuse--diffuse pairs (Li~3$s$--H~3$s$) carry weight
$1 \times 0.11 = 0.11$---an 81$\times$ discrimination.

The $R$-discrimination (ratio of correction at $R=1.5$ to $R=5.0$~bohr)
improves from 3.0$\times$ with uniform weighting to 9.3$\times$ with
Fock weighting.  For the Li~1$s$ orbital alone, the ratio improves from
19$\times$ to 74$\times$.

%% ------------------------------------------------------------------------
\subsection{Results of the Fock-Weighted Correction}
\label{sec:fock_results}
%% ------------------------------------------------------------------------

Table~\ref{tab:fock_scan} summarizes the $\lambda$-scan.

\begin{table}[h]
\centering
\begin{tabular}{lccc}
\toprule
$\lambda$ & Equilibrium? & $R_{\mathrm{eq}}$ (bohr) &
$D_e^{\mathrm{CP}}$ (Ha) \\
\midrule
0.00 & No & --- & --- \\
0.02 & No & --- & 0.091 at $R{=}3.0$ \\
0.05 & No & --- & uptick at $R{=}3.5$ \\
\textbf{0.10} & \textbf{Yes} & $\bm{\sim 3.8}$ &
\textbf{0.046} \\
0.20 & Yes & $\sim 4.4$ & 0.020 \\
0.50 & No (unbound) & --- & --- \\
\bottomrule
\end{tabular}
\caption{Fock-weighted kinetic correction scan.  At $\lambda = 0.1$,
a genuine potential energy minimum emerges at $R\approx 3.8$~bohr with
a repulsive wall below $R \approx 2.3$~bohr.}
\label{tab:fock_scan}
\end{table}

At $\lambda = 0.1$, the Fock-weighted correction creates a genuine
potential energy minimum with a proper repulsive wall: the molecule is
unbound below $R \approx 2.3$~bohr, maximally bound near
$R \approx 3.8$~bohr ($D_e^{\mathrm{CP}} \approx 0.046$~Ha), and
dissociates smoothly at large $R$.  This is a qualitative success---the
first equilibrium geometry produced by any balanced Hamiltonian
configuration---but the single-parameter model cannot simultaneously
reproduce both $R_{\mathrm{eq}}$ and $D_e$: at $\lambda = 0.02$,
the CP-corrected binding energy matches experiment
($D_e^{\mathrm{CP}} = 0.091$~Ha at $R=3.015$~bohr), but no equilibrium
forms.

%% ------------------------------------------------------------------------
\subsection{Remaining Limitations}
\label{sec:limitations}
%% ------------------------------------------------------------------------

The remaining discrepancy ($R_{\mathrm{eq}} = 3.8$ vs.\ experiment
$3.015$~bohr, $D_e = 0.046$ vs.\ $0.092$~Ha) likely reflects the
limitation of a purely diagonal correction.  In standard quantum
chemistry, kinetic repulsion includes both diagonal (orbital compression)
and off-diagonal (orthogonalization) contributions from the full kinetic
energy matrix $T_{AB} = \langle\phi_A|-\tfrac{1}{2}\nabla^2|\phi_B\rangle$.
The Fock-weighted correction captures only the diagonal component.

The BSSE of $0.115$~Ha \emph{strictly exceeds} the experimental
dissociation energy of $0.092$~Ha, confirming significant basis
deficiency at $n_{\max}=3$.  However, since the balanced PES has no
minimum, BSSE is not the mechanism for geometric contraction---it is a
separate issue of basis completeness that CP correction addresses for
vertical binding energies.

%% ========================================================================
\section{Conclusion}
\label{sec:conclusion}
%% ========================================================================

We have presented the first heteronuclear full configuration interaction
calculation on a discrete graph Hamiltonian.  The Topological LCAO
architecture concatenates atom-centered hydrogenic lattices via Mulliken
cross-nuclear attraction, Slater--Condon two-electron integrals, and
STO-overlap bridges, and diagonalizes the resulting Hamiltonian in the
full 367\,290-determinant Slater space using excitation-driven Direct CI.

For LiH at $n_{\max}=3$, the Boys--Bernardi counterpoise-corrected
vertical binding energy at $R=3.015$~bohr is
$D_e^{\mathrm{CP}}=0.110$~Ha, within 20\% of the experimental value
of $0.092$~Ha.  The balanced Hamiltonian produces a monotonically
attractive PES with no equilibrium geometry---a limitation traced to
the $R$-independence of the graph Laplacian's intra-atom kinetic
energy, not to basis truncation.

A 29-version diagnostic arc (v0.9.8--v0.9.37) systematically falsified
every competing hypothesis: basis set superposition error as a geometric
driver, integral incompleteness, orbital exponent relaxation, and the
SO(4) Sturmian Bond Sphere alternative.  The missing kinetic repulsion
was identified as a structural limitation of the LCAO concatenation,
and a Fock-weighted kinetic correction derived from Paper~7's
energy-shell constraint $p_0^2 = Z^2/n^2$ partially restores the
equilibrium ($R_{\mathrm{eq}} \approx 3.8$~bohr,
$D_e^{\mathrm{CP}} \approx 0.046$~Ha at $\lambda = 0.1$).

The $O(V)$ scaling of the underlying graph Hamiltonian is preserved
throughout, with the Direct CI solver achieving
$O(N_{\mathrm{SD}}\times N_{\mathrm{connected}})$ assembly cost.
The natural next steps are:
\begin{enumerate}
\item Extending the Fock-weighted correction to include off-diagonal
kinetic matrix elements between atomic centers, which would provide
the missing intermediate-$R$ binding;
\item Investigating whether the Bond Sphere geometry of Paper~8
naturally incorporates the required $R$-dependent kinetic coupling
through the SO(4) Wigner $D$-matrix;
\item Increasing $n_{\max}$ to assess basis convergence of the
vertical binding energy.
\item \textbf{Prolate spheroidal lattice:} The fundamental limitation
identified above---that the LCAO graph Laplacian provides no
$R$-dependent kinetic repulsion---arises because the molecular graph
is constructed as a concatenation of atomic $S^3$ subgraphs whose
adjacency matrices are fixed at construction time.  This limitation
is resolved by recognizing that the natural geometry for two-center
systems is the prolate spheroid, not the concatenated $S^3$.  In
prolate spheroidal coordinates $(\xi, \eta, \varphi)$, the two-center
Coulomb problem separates exactly, and discretizing the
Laplace-Beltrami operator on this manifold yields a molecular lattice
with $R$-dependence built into the metric.  For H$_2^+$, this approach
achieves $R_{\mathrm{eq}} = 2.001$~bohr (0.21\% error) and
$E_{\min} = -0.598$~Ha (0.70\% error) with zero fitted
parameters~\cite{loutey_paper11}.  The LCAO framework documented in
this paper remains valuable for understanding the anatomy of molecular
binding in a discrete framework, and the diagnostic arc
(Table~\ref{tab:versions}) provides a systematic record of approaches
that fail and why.
\end{enumerate}
The discovery that the Fock projection's energy-shell scale provides
the correct weighting for kinetic repulsion reinforces the physical
content of the conformal $S^3$ framework: the momentum parameter
$p_0^2 = Z^2/n^2$ is not merely a mathematical artifact of the
stereographic projection but encodes the ``information density'' of
each quantum state on the three-sphere.

%% ========================================================================
\appendix
\section{Version History}
\label{app:versions}
%% ========================================================================

Table~\ref{tab:versions} documents the complete diagnostic arc from
v0.9.11 (the counterpoise-corrected baseline) through v0.9.38, including
architecture, results, and root cause classification for each version.

\begin{table*}[t]
\centering
\small
\begin{tabular}{llccll}
\toprule
Version & Architecture & $D_e^{\mathrm{CP}}$ (Ha) & $R_{\mathrm{eq}}$ (bohr)
& Status & Root Cause \\
\midrule
v0.9.11 & LCAO baseline (balanced \texttt{exact+True})
  & 0.110 & No eq. & Baseline (20\%) & CP-corrected vertical $D_e$ \\
v0.9.12 & L\"owdin orthogonalization
  & --- & --- & Fail & Breaks variational principle \\
v0.9.13 & Mulliken exchange $K$
  & 0.110 & $\sim$2.5 & 19\% err & $K$ overcorrects $D_e$ \\
v0.9.15 & SO(4) Wigner $D$-matrix
  & --- & --- & Infrastructure & 44/44 tests, no production \\
v0.9.16 & $D$-matrix cross-nuclear (geom.\ mean)
  & --- & --- & Unbound & Coupling too weak \\
v0.9.17 & Shibuya--Wulfman off-diagonal
  & --- & --- & Unbound & SW coupling too weak \\
v0.9.18 & Hybrid (Fourier diag + SW off-diag)
  & 0.143 & $<$2.0 & 55\% err & Dissociation limit wrong \\
v0.9.20 & Sturmian diagonal
  & --- & --- & Unbound & 1$s$ diagonal positive \\
v0.9.21 & $p_0$-consistent $D$-matrix + self-loop
  & --- & --- & Unbound & Compactness overbinding \\
v0.9.22 & Atom-dependent $p_0$
  & 3.777 & $\sim$2.0 & $3.8\times$ overbind & Scales do not decouple \\
v0.9.23 & Shell-radius attenuation
  & --- & --- & Negligible & Correction $<$ 1 mHa \\
v0.9.24 & Energy decomposition diagnostic
  & --- & --- & Tool & Identifies Fourier diagonal \\
v0.9.25 & Ohno--Klopman cross-atom $V_{ee}$
  & --- & --- & Negligible & Cross-atom $J$ irrelevant \\
v0.9.26 & Angular-weighted $V_{ee}$
  & --- & --- & Negligible & $\Delta E \sim 0.4$ mHa \\
v0.9.27 & SO(4) CG cross-center $V_{ee}$
  & --- & --- & Unbound & $J_{\mathrm{cross}} > J_{\mathrm{same}}$ \\
v0.9.28 & Bond sphere node weights
  & --- & --- & Unbound & $\beta_H$ 46\% short \\
v0.9.29 & Prolate spheroidal MO betas
  & --- & --- & Unbound & LiH mapping ad hoc \\
v0.9.30 & MO-to-atom projection
  & --- & --- & Unphysical & $E=-42.5$ Ha ($5.3\times$) \\
v0.9.31 & MO Sturmian FCI (full)
  & --- & --- & Unbound/overbinding & $\mathbf{H}\!\propto\!\mathbf{S}$, eigenvalues $\gamma$-independent \\
v0.9.32 & MO Sturmian FCI (variant)
  & --- & --- & Unbound/overbinding & $\mathbf{H}\!\propto\!\mathbf{S}$, eigenvalues $\gamma$-independent \\
v0.9.33 & MO Sturmian FCI (variant)
  & --- & --- & Unbound/overbinding & $\mathbf{H}\!\propto\!\mathbf{S}$, eigenvalues $\gamma$-independent \\
v0.9.34 & Dual-$p_0$ extension
  & --- & --- & Unbound/overbinding & $\mathbf{H}\!\propto\!\mathbf{S}$, eigenvalues $\gamma$-independent \\
v0.9.35 & All-$l$ cross-atom $V_{ee}$
  & 0.110 & No eq. & $\Delta E = 0.4$~mHa & Negligible correction \\
v0.9.36 & Orbital exponent relaxation
  & 0.433 & No eq. & $4.7\times$ overbind & Uniform PES shift \\
v0.9.37 & Hamiltonian balance diagnostic
  & 0.110 & No eq. & Corrected & Imbalance identified \\
v0.9.38 & Numba-accelerated assembly
  & 0.110 & No eq. & $738\times$ speedup & Python loop overhead eliminated \\
\bottomrule
\end{tabular}
\caption{Diagnostic arc version history (v0.9.11--v0.9.38).
The LCAO baseline (v0.9.11, balanced \texttt{exact+True}) achieves
correct-order-of-magnitude binding ($D_e^{\mathrm{CP}} = 0.110$~Ha,
20\% error) but lacks equilibrium geometry due to missing kinetic
repulsion.}
\label{tab:versions}
\end{table*}

%% ========================================================================
%% REFERENCES
%% ========================================================================

\begin{thebibliography}{9}

\bibitem{paper0}
J.~Loutey, ``The Geometric Atom: Quantum State Space as a Packing Problem,''
GeoVac Paper~0 (2026).

\bibitem{paper1}
J.~Loutey, ``The Geometric Atom: Quantum Mechanics as a Packing Problem,''
GeoVac Paper~1 (2026).

\bibitem{paper7}
J.~Loutey, ``The Dimensionless Vacuum: Scale-Invariant Graph Laplacian
on $S^3$,'' GeoVac Paper~7 (2026).

\bibitem{paper8}
J.~Loutey, ``Bond Sphere Theory and Sturmian Selection Rules,''
GeoVac Papers~8--9 (2026).

\bibitem{fock1935}
V.~Fock, ``Zur Theorie des Wasserstoffatoms,''
Z.~Phys. \textbf{98}, 145 (1935).

\bibitem{boys_bernardi}
S.~F.~Boys and F.~Bernardi, ``The calculation of small molecular
interactions by the differences of separate total energies.  Some
procedures with reduced errors,''
Mol.~Phys. \textbf{19}, 553 (1970).

\bibitem{shull_lowdin}
H.~Shull and P.-O.~L\"owdin, ``Superposition of Configurations and
Natural Spin Orbitals.\ Applications to the He Problem,''
J.~Chem.~Phys. \textbf{25}, 1035 (1956).

\bibitem{huber}
K.~P.~Huber and G.~Herzberg, \emph{Molecular Spectra and Molecular
Structure IV: Constants of Diatomic Molecules}
(Van Nostrand Reinhold, New York, 1979).

\bibitem{knowles_handy}
P.~J.~Knowles and N.~C.~Handy, ``A new determinant-based full
configuration interaction method,''
Chem.~Phys.~Lett. \textbf{111}, 315 (1984).

\bibitem{numba}
S.~K.~Lam, A.~Pitrou, and S.~Seibert, ``Numba: A LLVM-based Python
JIT compiler,'' in \emph{Proceedings of the Second Workshop on the LLVM
Compiler Infrastructure in HPC} (ACM, 2015), pp.~1--6.

\bibitem{loutey_paper11}
J.~Loutey,
``The Molecular Fock Projection: Diatomic Molecules as Prolate
Spheroidal Lattices,''
GeoVac Paper~11 (2026).

\end{thebibliography}

\end{document}
